\chapter{激光冷却、Ramsey 干涉、原子钟与磁阱}
\label{ch:laser-cooling-clock}

\section{散射力、摩擦与反冲扩散}

冷却操纵的是原子质心自由度。下文用 $M$ 表示\emph{整颗原子}的质量；电子质量 $m_e$ 不出现在质心 Fokker--Planck 方程中。内部两能级稳态则直接继承\eqnrefs{eq:two-level-saturation}。

原子吸收一个沿 \(\mathbf k\) 方向的光子，动量改变

\[
\Delta\mathbf p=\hbar\mathbf k.
\]

自发辐射方向随机，长时间平均动量接近零。因此平均散射力

\[
\mathbf F_{sc}=\hbar\mathbf k\,R_{sc}
\]

其中散射率

\[
R_{sc}=\Gamma_{\rm sp}\rho_{ee}^{(\rm ss)}.
\]

于是

\[
\mathbf F_{sc}
=\hbar\mathbf k\frac{\Gamma_{\rm sp}}2
\frac{s_0}{1+s_0+[2(\Delta-\mathbf k\cdot\mathbf v)/\Gamma_{\rm sp}]^2}
.
\]

Doppler 位移通过

\[
\Delta\rightarrow\Delta-\mathbf k\cdot\mathbf v
\]

进入。

取两束等强、相向传播的红失谐激光。对低饱和 \(s_0\ll1\)，两束光的饱和效应可以在首阶分别处理。沿 \(x\) 方向的总力为

\begin{proof}[解]
记 Lorentz 线型
\begin{equation*}
f(\delta)=\frac{1}{1+(2\delta/\Gamma_{\rm sp})^2}.
\end{equation*}
沿 $+x$ 与 $-x$ 两束等强红失谐光，原子所受单束散射力分别为
\begin{equation*}
F_\pm(v)=\frac{\hbar k\Gamma_{\rm sp}s_0}{2}\,f(\Delta\mp kv),
\end{equation*}
其中 $\mp kv$ 来自 Doppler 失谐 $\Delta\mp kv$。总力为
\begin{equation*}
F(v)=F_+(v)-F_-(v)
=\frac{\hbar k\Gamma_{\rm sp}s_0}{2}\bigl[f(\Delta-kv)-f(\Delta+kv)\bigr].
\end{equation*}
在 $kv\ll\Gamma_{\rm sp}$ 的低速区对 $kv$ 作一阶 Taylor 展开，
\begin{equation*}
f(\Delta\mp kv)=f(\Delta)\mp kv\,f'(\Delta)+\frac{(kv)^2}{2}\,f''(\Delta)+O(v^3).
\end{equation*}
代入总力，偶次项相消而奇次项相加，
\begin{equation*}
F(v)=\frac{\hbar k\Gamma_{\rm sp}s_0}{2}\bigl[-2kv\,f'(\Delta)+O(v^3)\bigr]
=-\hbar k^2\Gamma_{\rm sp}s_0\,f'(\Delta)\,v+O(v^3).
\end{equation*}
为求 $f'(\Delta)$，把 $f$ 写成 $f(\delta)=[1+u(\delta)]^{-1}$，$u(\delta)=(2\delta/\Gamma_{\rm sp})^2$，于是
\begin{equation*}
f'(\delta)=-\frac{u'(\delta)}{[1+u(\delta)]^2},
\qquad
u'(\delta)=\frac{8\delta}{\Gamma_{\rm sp}^2},
\end{equation*}
因此
\begin{equation*}
f'(\Delta)
=-\frac{8\Delta/\Gamma_{\rm sp}^2}{\bigl[1+(2\Delta/\Gamma_{\rm sp})^2\bigr]^2}.
\end{equation*}
代回 $F(v)$ 并整理，
\begin{equation*}
F(v)=-\eta\,v+O(v^3),
\qquad
\eta=-\frac{8\hbar k^2 s_0\,(\Delta/\Gamma_{\rm sp})}{\bigl[1+(2\Delta/\Gamma_{\rm sp})^2\bigr]^2}.
\end{equation*}
故阻尼系数 $\eta$ 的符号由谱线斜率 $f'(\Delta)$ 决定：红失谐 $\Delta<0$ 给出 $f'(\Delta)>0$，于是 $\eta>0$，恢复性摩擦；蓝失谐反向。即正文的 $F=-\eta v$。
\end{proof}

\[
F(v)=\frac{\hbar k\Gamma_{\rm sp}}{2}s_0
\left[
\frac{1}{1+[2(\Delta-kv)/\Gamma_{\rm sp}]^2}
-\frac{1}{1+[2(\Delta+kv)/\Gamma_{\rm sp}]^2}
\right].
\]

若

\[
|kv|\ll\Gamma_{\rm sp},|\Delta|,
\]

对速度展开到一阶，常数项因两束光对称而相消，得到

\[
F(v)=-\eta v+O(v^3)
\]

其中

\[
\eta
=-\frac{8\hbar k^2s_0(\Delta/\Gamma_{\rm sp})}
{[1+(2\Delta/\Gamma_{\rm sp})^2]^2}.
\]

红失谐 \(\Delta<0\) 时 \(\eta>0\)，所以速度受到线性阻尼。蓝失谐会改变符号并导致反阻尼。

散射力虽然提供阻尼，每次自发辐射却带来随机反冲。低饱和近似下，可以把速度分布的慢演化写成 Fokker--Planck 方程

\[
\partial_t W(p,t)
=\frac{\eta}{M}\partial_p[pW]
+D_p\partial_p^2W,
\]

其中 \(D_p\) 是动量扩散系数。稳态为高斯，满足 Einstein 型关系

\[
k_BT=\frac{D_p}{\eta}.
\]

本讲义采用 Fokker--Planck 形式

\[
\partial_tW=-\partial_p(AW)+D_p\partial_p^2W,
\]

所以随机过程满足

\[
\frac{\dd}{\dd t}\langle(\Delta p)^2\rangle=2D_p.
\]

在最简单的一维反冲模型中，每次散射事件包含一次定向吸收和一次随机自发辐射。两束光在 \(v=0\) 附近的总散射率为

\[
R_{\rm tot}
=\frac{\Gamma_{\rm sp}s_0}
{1+(2\Delta/\Gamma_{\rm sp})^2}
+O(s_0^2).
\]

若把自发辐射也简化成沿 \(\pm x\) 等概率反冲，则单次完整散射循环对 \(p_x\) 方差的平均贡献是 \(2(\hbar k)^2\)。Poisson 统计给出

\[
\frac{\dd}{\dd t}\langle(\Delta p_x)^2\rangle
=2(\hbar k)^2R_{\rm tot},
\]

因此

\begin{equation}
D_p
=(\hbar k)^2
\frac{\Gamma_{\rm sp}s_0}
{1+(2\Delta/\Gamma_{\rm sp})^2}
\label{eq:doppler-diffusion}
\end{equation}

在三维真实偶极辐射中，自发发射的角分布会改变数值系数，但不会改变 \(D_p\propto(\hbar k)^2R_{\rm sc}\) 的结构，也不改变 Doppler 极限来自“摩擦与反冲扩散平衡”的物理机制。

把吸收与自发辐射的反冲扩散和上一节的摩擦系数代入，在一维、低饱和二能级模型中得到

\[
k_BT(\Delta)
=\frac{\hbar\Gamma_{\rm sp}}{4}
\frac{1+(2\Delta/\Gamma_{\rm sp})^2}
{|2\Delta|/\Gamma_{\rm sp}}.
\]

确实，把 \eqnrefs{eq:doppler-diffusion} 与上一节

\[
\eta
=-\frac{8\hbar k^2s_0(\Delta/\Gamma_{\rm sp})}
{[1+(2\Delta/\Gamma_{\rm sp})^2]^2}
\]

代入 \(k_BT=D_p/\eta\)，所有 \(k^2\) 和低饱和强度因子 \(s_0\) 消去，便得到上式。温度在低饱和极限与激光强度无关，正是因为增大强度会同比增大阻尼和反冲噪声。

令 \(x=|2\Delta|/\Gamma_{\rm sp}\)，右端正比于 \(x+x^{-1}\)，所以极小值在 \(x=1\)，即

\[
\Delta=-\frac{\Gamma_{\rm sp}}{2}.
\]

由此得到 Doppler 极限

\[
T_D=\frac{\hbar\Gamma_{\rm sp}}{2k_B}.
\]

温度关于失谐的无量纲曲线以及最优红失谐位置见\figref{fig:doppler-limit}。

\begin{proof}[解]
在时间 $\dd t$ 内发生一次散射的概率为 $R_{\rm tot}\dd t+O(\dd t^2)$，散射事件近似 Poisson 过程。记 $W(p,t)$ 为纵向动量分布，单步跳变 $p\to p+\delta p$ 的主方程为
\begin{equation*}
\partial_t W(p,t)
= R_{\rm tot}\!\int\! \dd(\delta p)\,
 \bigl[W(p-\delta p,t)-W(p,t)\bigr]\,P(\delta p),
\end{equation*}
其中 $P(\delta p)$ 为单次反冲分布。对 $W(p-\delta p,t)$ 在 $p$ 处 Taylor 展开，
\begin{equation*}
W(p-\delta p,t)
= W - \delta p\,\partial_p W
 + \frac{(\delta p)^2}{2}\,\partial_p^2 W + O\bigl((\delta p)^3\bigr).
\end{equation*}
代回主方程并按矩 $\langle\delta p\rangle=0$、$\langle(\delta p)^2\rangle=2(\hbar k)^2$ 截断到二阶（Kramers--Moyal），因此
\begin{equation*}
\partial_t W
= -\partial_p(AW)+\frac12\,\partial_p^2(BW),
\qquad
A=\langle\delta p\rangle R_{\rm tot}=0,
\quad
B=\langle(\delta p)^2\rangle R_{\rm tot}=2(\hbar k)^2 R_{\rm tot}.
\end{equation*}
扩散系数 $D_p=B/2=(\hbar k)^2 R_{\rm tot}$，即\eqnrefs{eq:doppler-diffusion}。低速摩擦 $F=-\eta v=-\eta p/M$ 给出漂移 $A(p)=-\eta p/M$，于是 Fokker--Planck 方程化为
\begin{equation*}
\partial_t W
= \frac{\eta}{M}\,\partial_p(pW)+D_p\,\partial_p^2 W.
\end{equation*}
稳态零概率流 $J=(\eta/M)pW+D_p\,\partial_p W=0$，分离变量，
\begin{equation*}
\frac{\partial_p W}{W}
= -\frac{\eta}{M D_p}\,p,
\end{equation*}
积分得
\begin{equation*}
W_{\rm ss}(p)\propto
\exp\!\left[-\frac{\eta\,p^2}{2 M D_p}\right]
=\exp\!\left[-\frac{p^2}{2 M k_B T}\right],
\qquad
k_B T=\frac{D_p}{\eta}.
\end{equation*}
代入 $D_p=(\hbar k)^2 R_{\rm tot}$ 与 $\eta=-8\hbar k^2 s_0(\Delta/\Gamma_{\rm sp})/[1+(2\Delta/\Gamma_{\rm sp})^2]^2$，$k^2$ 与 $s_0$ 消去，记 $x=|2\Delta|/\Gamma_{\rm sp}$，整理得
\begin{equation*}
k_B T
= \frac{\hbar\Gamma_{\rm sp}}{4}\left(x+\frac1x\right).
\end{equation*}
由 AM--GM 不等式 $x+x^{-1}\ge2$，等号在 $x=1$ 取得，因此最优失谐 $\Delta=-\Gamma_{\rm sp}/2$（红失谐侧），于是
\begin{equation*}
k_B T_D=\frac{\hbar\Gamma_{\rm sp}}{2},
\qquad
T_D=\frac{\hbar\Gamma_{\rm sp}}{2 k_B}.
\end{equation*}
\end{proof}

\begin{figure}[tbp]
\centering
\includegraphics[width=.83\textwidth]{doppler_limit.pdf}
\caption{低饱和二能级 Doppler 模型中的平衡温度随失谐的变化。极小值出现在\(\Delta=-\Gamma_{\rm sp}/2\)。}
\label{fig:doppler-limit}
\end{figure}

这个极限依赖``二能级、低饱和、仅 Doppler 摩擦与单光子反冲扩散''的模型。存在 Zeeman 简并和偏振梯度时，额外的亚 Doppler 机制可以把温度继续降到 \(T_D\) 以下。

\section{内部简并与亚 Doppler 冷却}
Doppler 极限并不是绝对下限。具有简并 Zeeman 子能级的原子在偏振梯度场中会发生位置依赖的 light 位移。

例如两束正交线偏振光反向传播：

\[
\mathbf E^{(+)}(z)
=\frac{E_0}{2}
(\hat x \ee^{\ii kz}+\hat y \ee^{-\ii kz}).
\]

重新选取正交线偏振基，可看出偏振随 \(z\) 周期变化：

\[
\text{linear}\rightarrow \sigma^-\rightarrow\text{linear}\rightarrow\sigma^+\rightarrow\cdots
\]

不同 \(m_F\) 子态由于 CG 系数不同，得到不同的 AC Stark 势 \(U_{m_F}(z)\)。

以最简单的 \(J_g=1/2\to J_e=3/2\) 模型为例。对基态 \(m_g=+1/2\)，\(\sigma^+\) 跃迁到 \(m_e=+3/2\) 的 CG 强度为 1，而 \(\sigma^-\) 到 \(m_e=-1/2\) 的强度为 \(1/3\)；对 \(m_g=-1/2\) 两者交换。因此在远失谐、低饱和条件下，若局域圆偏振分量强度为 \(I_\pm(z)\)，两条基态 light 位移具有结构

\[
U_+(z)=C\left[I_+(z)+\frac13 I_-(z)\right],
\qquad
U_-(z)=C\left[I_-(z)+\frac13 I_+(z)\right],
\]

其中 \(C\propto1/\Delta\) 包含共同的偶极矩和失谐因子。lin\(\perp\)lin 场中 \(I_+\) 与 \(I_-\) 在空间上错相振荡，因此 \(U_+\) 与 \(U_-\) 也是相互错位的周期势。光抽运率又在局域圆偏振最强处最大，恰好使原子更容易在“爬到某一势能山顶”附近跳到另一条较低势曲线。Sisyphus 图像由此从“偏振变化”定量落到“CG 加权的两条位置依赖光移”。

原子沿势坡上爬，动能转化为势能；在接近势顶时被光学 pumping 到另一个更低的子态，再次开始爬坡。重复这一过程就像希腊神话中的 Sisyphus 不断推石上山，因此称 Sisyphus cooling。

原子电子总角动量

\[
\mathbf J=\mathbf L+\mathbf S,
\]

与核自旋 \(\mathbf I\) 耦合形成

\[
\mathbf F=\mathbf I+\mathbf J.
\]

磁偶极超精细哈密顿量

\[
H_{hf}=A_{\rm hf}\mathbf I\cdot\mathbf J.
\]

利用

\[
\mathbf I\cdot\mathbf J
=\frac12(\mathbf F^2-\mathbf I^2-\mathbf J^2),
\]

得到

\[
E_F
=\frac{A_{\rm hf}}2
[F(F+1)-I(I+1)-J(J+1)]
.
\]

例如 \(^{133}\mathrm{Cs}\) 基态 \(6S_{1/2}\)：

\[
I=\frac72,\qquad J=\frac12,
\]

所以

\[
F=3,4.
\]

这两个超精细态之间的跃迁频率定义了 SI 秒的基础。

\section{Ramsey 干涉、频率基准与相对论修正}
Ramsey 方法把``连续长时间驱动''改写为两个短 \(\pi/2\) 脉冲，中间留出自由相位积累时间 \(T\)。在理想瞬时脉冲近似下，取第一脉冲相位为零，\(\pi/2\) 操作可写为

\[
U_{\pi/2}=\frac1{\sqrt2}
\begin{pmatrix}
1&-\ii\\
-\ii&1
\end{pmatrix}.
\]

旋转表象中自由演化传播子为

\[
U_0(T)=
\begin{pmatrix}
 \ee^{+\ii\Delta T/2}&0\\
0&\ee^{-\ii\Delta T/2}
\end{pmatrix}.
\]

从基态 \((1,0)^T\) 出发，完整序列

\[
|\psi_f\rangle
=U_{\pi/2}U_0(T)U_{\pi/2}|g\rangle
\]

给出

\[
P_e(T)=\frac12[1+\cos(\Delta T)]
\]

（改变第二脉冲相位只会平移 fringe）。因此中央条纹的频率尺度由 \(1/T\) 决定，而不是由短脉冲宽度决定。

\begin{proof}[解]
记基态与激发态列向量 $|g\rangle=(1,0)^T$、$|e\rangle=(0,1)^T$。第一束同相位 $\pi/2$ 脉冲的 SU(2) 作用为
\begin{equation*}
U_{\pi/2}=\frac1{\sqrt2}
\begin{pmatrix}1&-\ii\\-\ii&1\end{pmatrix},
\end{equation*}
作用在 $|g\rangle$ 上，
\begin{equation*}
U_{\pi/2}|g\rangle
=\frac1{\sqrt2}
\begin{pmatrix}1&-\ii\\-\ii&1\end{pmatrix}
\begin{pmatrix}1\\0\end{pmatrix}
=\frac1{\sqrt2}\begin{pmatrix}1\\-\ii\end{pmatrix}.
\end{equation*}
在旋转表象中，自由演化对角传播子为
\begin{equation*}
U_0(T)=
\begin{pmatrix}\ee^{+\ii\Delta T/2}&0\\0&\ee^{-\ii\Delta T/2}\end{pmatrix},
\end{equation*}
作用后
\begin{equation*}
U_0(T)U_{\pi/2}|g\rangle
=\frac1{\sqrt2}
\begin{pmatrix}\ee^{+\ii\Delta T/2}&0\\0&\ee^{-\ii\Delta T/2}\end{pmatrix}
\begin{pmatrix}1\\-\ii\end{pmatrix}
=\frac1{\sqrt2}
\begin{pmatrix}\ee^{+\ii\Delta T/2}\\-\ii\ee^{-\ii\Delta T/2}\end{pmatrix}.
\end{equation*}
再作用第二束同相位 $\pi/2$ 脉冲，
\begin{align*}
|\psi_f\rangle
&=U_{\pi/2}\,U_0(T)\,U_{\pi/2}|g\rangle \\
&=\frac1{\sqrt2}
\begin{pmatrix}1&-\ii\\-\ii&1\end{pmatrix}
\cdot\frac1{\sqrt2}
\begin{pmatrix}\ee^{+\ii\Delta T/2}\\-\ii\ee^{-\ii\Delta T/2}\end{pmatrix} \\
&=\frac12
\begin{pmatrix}
\ee^{+\ii\Delta T/2}-\ee^{-\ii\Delta T/2}\\
-\ii\bigl(\ee^{+\ii\Delta T/2}+\ee^{-\ii\Delta T/2}\bigr)
\end{pmatrix}.
\end{align*}
由 Euler 公式 $\ee^{\ii x}+\ee^{-\ii x}=2\cos x$、$\ee^{\ii x}-\ee^{-\ii x}=2\ii\sin x$，于是
\begin{equation*}
|\psi_f\rangle
=\begin{pmatrix}
\ii\sin(\Delta T/2)\\
-\ii\cos(\Delta T/2)
\end{pmatrix},
\end{equation*}
激发态振幅
\begin{equation*}
c_e=-\ii\cos\frac{\Delta T}{2}.
\end{equation*}
取模平方，
\begin{equation*}
P_e=|c_e|^2=\cos^2\frac{\Delta T}{2}
=\frac12\bigl[1+\cos(\Delta T)\bigr].
\end{equation*}
自由演化时间 $T$ 把微小失谐放大为相位 $\Delta T$，这正是 Ramsey spectroscopy 的分辨率由 $1/T$ 而非脉冲宽度决定的来源。
\end{proof}

若一次测量包含 \(N\) 个互不纠缠原子，在最敏感工作点上二项分布投影噪声给出 \(\Delta P\sim1/(2\sqrt N)\)，而 \(|\partial P/\partial\phi|=1/2\)，所以

\[
\Delta\phi\sim\frac1{\sqrt N},
\qquad
\Delta\omega\sim\frac1{T\sqrt N}.
\]

这就是标准量子极限的尺度。有限脉冲时间、局域振荡器噪声、原子数涨落和 decoherence 会改变实际 Allan deviation，但不会改变这里展示的基本干涉机制。

频率基准一旦稳定到足够精度，钟差本身就成为可测距离。

GPS 接收机要解四个未知量：

\[
(x,y,z,\delta t).
\]

每颗卫星给出一个伪距方程

\[
\rho_i
=c(t_{rec}-t_i)
=\sqrt{(x-x_i)^2+(y-y_i)^2+(z-z_i)^2}+c\delta t.
\]

因此至少需要四颗卫星。

时间误差 \(\delta t\) 会直接转化为空间误差

\[
\delta r\sim c\delta t.
\]

例如

\[
1\,\text{ns}\Rightarrow 0.30\,\text{m}.
\]

所以卫星导航本质上是一项高精度时间测量技术。

弱场极限中，不同重力势 \(\Phi\) 下时钟频率满足

\[
\frac{\Delta\nu}{\nu}
\approx\frac{\Delta\Phi}{c^2}.
\]

在均匀重力场中，高度差 \(H\)：

\[
\Delta\Phi=gH,
\]

所以

\[
\frac{\Delta\nu}{\nu}\approx\frac{gH}{c^2}
.
\]

这既可从广义相对论理解，也可借助等效原理在加速电梯中用 Doppler 位移得到。

\section{磁俘获、绝热内部态与合成规范场}
Zeeman 相互作用

\[
H_Z=-\boldsymbol\mu\cdot\mathbf B.
\]

绝热情况下，某个内部态的能量近似为

\[
U(\mathbf r)=\mu_{eff}|\mathbf B(\mathbf r)|.
\]

若 \(\mu_{eff}>0\)，能量在弱磁场处更低，称 low-场-seeking 态，可被 \(|B|\) 的极小值束缚。

在无电流区域

\[
\nabla\cdot\mathbf B=0,
\qquad
\nabla\times\mathbf B=0,
\]

因此各分量满足 Laplace 方程。静态自由空间磁场不能产生稳定的 \(|B|\) 局域极大值，因此 high-场-seeking 态不能通过普通静态磁场极大值稳定俘获。

这个结论可以直接从 \(|\mathbf B|\) 的拉普拉斯算子算出来。记 \(B=|\mathbf B|\)，则

\[
\partial_iB=\frac{B_j\partial_iB_j}{B},
\]

进一步

\[
\partial_i^2B
=\frac{(\partial_iB_j)^2+B_j\partial_i^2B_j}{B}
-\frac{(B_j\partial_iB_j)^2}{B^3}.
\]

对空间指标 \(i\) 求和，并在无源区使用 \(\nabla^2B_j=0\)：

\[
\nabla^2B
=\frac{\sum_{ij}(\partial_iB_j)^2}{B}
-\frac{\sum_i(B_j\partial_iB_j)^2}{B^3}.
\]

Cauchy--Schwarz 不等式给出

\[
(B_j\partial_iB_j)^2
\le B^2\sum_j(\partial_iB_j)^2,
\]

因此

\begin{equation}
\nabla^2|\mathbf B|\ge0
\label{eq:magnetic-field-subharmonic}
\end{equation}

只要 \(B\neq0\)，\(|\mathbf B|\) 是 subharmonic 函数，因而不能在无源区有严格局域极大值。这就是静态磁阱版的 Earnshaw 限制。极小值则不被禁止，所以 low-场 seeker 可以被 Ioffe--Pritchard 等非零极小值磁阱稳定束缚；保持非零偏置场还能避免零场处的 Majorana 自旋-flip 损失。

这不是说``超精细态不能被磁阱束缚''，相反，磁阱实验中正是选择合适的低场趋向超精细态。

若内部哈密顿量随位置变化：

\[
H_{int}(\mathbf r)|\chi_n(\mathbf r)\rangle
=\varepsilon_n(\mathbf r)|\chi_n(\mathbf r)\rangle,
\]

展开总波函数

\[
|\Psi\rangle
=\sum_n\psi_n(\mathbf r)|\chi_n(\mathbf r)\rangle.
\]

动量作用时会打到位置依赖的内部基，产生 Berry 联络：

\[
\mathbf A_{mn}
=\ii\hbar\langle\chi_m|\nabla\chi_n\rangle
.
\]

绝热投影到单个能带后

\[
H_{eff}
=\frac1{2M}[\mathbf p-\mathbf A_n(\mathbf r)]^2
+\varepsilon_n(\mathbf r)+\Phi_n(\mathbf r).
\]

其中单带几何标量势来自被绝热消去的其他内部态：

\[
\Phi_n(\mathbf r)
=\frac{\hbar^2}{2M}
\sum_{m\neq n}
|\langle\chi_m|\nabla\chi_n\rangle|^2.
\]

Berry 曲率则为

\[
\mathbf B_{\rm syn}=\nabla\times\mathbf A_n.
\]

因此“合成磁场”并不是把真实磁场直接替换成另一个实验旋钮，而是原子质心在位置依赖内部本征基中绝热运动时积累的几何相位所产生的有效 Lorentz 结构。这与\chapref{ch:born-huang}中的分子核几何规范势是同一个数学机制。

这就是冷原子中 synthetic 规范场的基本来源。

\begin{proof}[解]
以 \(J_g=1/2\to J_e=3/2\) 的 lin\(\perp\)lin 简化模型为例（\(U_0\) 为最大光位移深度，\(\Omega\) 为共振 Rabi 频率，\(\Gamma=\Gamma_{\rm sp}\) 为自然自发辐射线宽，\(\Delta\) 为光失谐），两条绝热光位移势和代表性抽运率写成
\begin{align}
 U_\pm(z)&=\frac{U_0}{2}[1\pm\cos(2kz)],
 \label{eq:sisyphus-light-shift}\\
 \Gamma_{\rm pump}(z)&=\Gamma'\sin^2(2kz),\qquad
 \Gamma'\simeq\frac{\Gamma\Omega^2}{4\Delta^2}
 \quad(|\Delta|\gg\Gamma).
 \label{eq:sisyphus-pump-rate}
\end{align}
式 \eqnrefs{eq:sisyphus-light-shift} 与式 \eqnrefs{eq:sisyphus-pump-rate} 中，具体 Clebsch--Gordan 系数会改变数值前因子，但不改变下面的尺度关系。

令 \(p_\pm(z)\) 为原子沿 \(z=vt\) 运动时的局域内态布居。它们满足速率方程
\[
 v\partial_z
 \begin{pmatrix}p_+\\p_-\end{pmatrix}
 =\mathsf W(z)
 \begin{pmatrix}p_+\\p_-\end{pmatrix},
 \qquad p_++p_-=1,
\]
其中 \(\mathsf W\) 由局域 \(\sigma^\pm\) 分量的抽运率决定。在
\(kv\ll\Gamma'\) 时作绝热响应展开
\[
 \vb p(z)=\vb p^{(0)}(z)+v\,\vb p^{(1)}(z)+O(v^2),\qquad
 \mathsf W\vb p^{(1)}=\partial_z\vb p^{(0)} .
\]
周期平均偶极力
\[
 \overline F=-\frac{k}{\pi}\int_0^{\pi/k}
 \sum_{\sigma=\pm}p_\sigma(z)\,\partial_zU_\sigma(z)\dd z
\]
的零阶项为零，一阶项给出 \(\overline F=-\alpha_{\rm Sis}v+O(v^2)\)，其中
\begin{equation}
 \alpha_{\rm Sis}
 =\frac{k}{\pi}\int_0^{\pi/k}
 \sum_{\sigma=\pm}p_\sigma^{(1)}(z)\,\partial_zU_\sigma(z)\dd z
 =C_\alpha\,\frac{k^2|U_0|}{\Gamma'} .
 \label{eq:sisyphus-friction}
\end{equation}
\(C_\alpha>0\) 是由跃迁角动量和偏振几何决定的无量纲常数。这个结果来自布居对移动光势的滞后；不能用“每个周期固定损失 \(U_0\)”的常功率估计替代线性响应。

动量扩散同时来自自发发射反冲和随机切换光势造成的偶极力涨落：
\[
 D_p=C_{\rm sp}(\hbar k)^2\Gamma'
 +C_{\rm dip}\frac{k^2U_0^2}{\Gamma'}.
\]
这些无量纲系数由发射角分布和内态结构决定。在典型偏振梯度冷却区
\(|U_0|\gg\hbar\Gamma'\)，偶极力涨落项主导。Fokker--Planck 稳态的
\(k_BT=D_p/\alpha_{\rm Sis}\) 此时给出
\begin{equation}
 k_BT_{\rm Sis}=C_T|U_0|,\qquad
 C_T=\frac{C_{\rm dip}}{C_\alpha}=O(1).
 \label{eq:sisyphus-temperature}
\end{equation}
所以亚 Doppler 标度由可调光位移而不是自然线宽控制；数值系数必须由完整多能级光学 Bloch 方程或实验几何确定。继续降低 \(|U_0|\) 最终会触及反冲能、非绝热跃迁和俘获效率等限制，不能由上式推出任意低温。

作为对照，二能级 Doppler 模型在弱饱和、最优失谐
\(\Delta=-\Gamma_{\rm sp}/2\) 下给出
\begin{equation}
 k_BT_D=\frac{\hbar\Gamma_{\rm sp}}{2}.
 \label{eq:phys-lc-doppler-limit}
\end{equation}
偏振梯度冷却能够低于这一数值，正因为其摩擦来自内态相关的光位移，而扩散仍由随机反冲决定。
\end{proof}
